% GENERATED, edit the French unit or the translation table
\mxnewpage\section*{EV3-T2 -- Mock exam 2}\label{EV3-T2}
\addcontentsline{toc}{section}{EV3-T2 -- Mock exam 2}
\textbf{Practice paper}. Programme studied in chapters 1 to 10.\par Duration: 2 h. No calculator. Marked out of 20. 1--3 April 2027.
\subsection*{Part 1 -- Fluency (6 points)}
Twelve questions, \num{0,5} point each. Only one correct answer; no justification required. An incorrect or missing answer scores 0, with no loss of points.
\begin{enumerate}
\item Calculate $2+2\times3$.\par A. $11$\quad B. $7$\quad C. $12$\quad D. $8$
\item Calculate $\frac12+\frac13$.\par A. $\frac{1}{6}$\quad B. $\frac{2}{5}$\quad C. $\frac{5}{3}$\quad D. $\frac{5}{6}$
\item Calculate $\dfrac{2^2}{2^{0}}$.\par A. $7$\quad B. $4$\quad C. $2$\quad D. $1$
\item A value of 50 increases by \pc{20}. Final value?\par A. $40$\quad B. $70$\quad C. $10$\quad D. $60$
\item A value of 80 is obtained after a \pc{20} decrease. Initial value?\par A. $96$\quad B. $60$\quad C. $100$\quad D. $64$
\item A \pc{25} increase then a \pc{20} decrease: overall coefficient?\par A. $\frac{4}{5}$\quad B. $\frac{21}{20}$\quad C. $1$\quad D. $\frac{9}{10}$
\item The solution of $2x+2=0$ is\par A. $-1$\quad B. $4$\quad C. $-4$\quad D. $1$
\item The slope between $\pt{0;2}$ and $\pt{2;8}$ is\par A. $3$\quad B. $2$\quad C. $6$\quad D. $-3$
\item The mean of $2, 4, 6$ is\par A. $2$\quad B. $4$\quad C. $5$\quad D. $12$
\item A sector of $90^\circ$ represents what fraction of the total?\par A. $\frac{1}{4}$\quad B. $\frac{1}{2}$\quad C. $\frac{1}{3}$\quad D. $\frac{3}{4}$
\item A probability is $\dfrac{2}{5}$. Its complement is\par A. $\frac{5}{2}$\quad B. $\frac{3}{5}$\quad C. $\frac{1}{5}$\quad D. $\frac{2}{5}$
\item For $f(x)=x^2+2$, $f(-2)$ is\par A. $6$\quad B. $-2$\quad C. $4$\quad D. $8$
\end{enumerate}
\mxnewpage\subsection*{Part 2 -- Independent exercises (14 points)}
Every step must be justified. Exact results are expected unless stated otherwise.
\mxneed{12}
\subsection*{Exercise 1 -- Expectation, variance and gain (5 points)}
The random variable $X$ takes the values 0, 2 and 6 with respective probabilities $\dfrac12$, $\dfrac14$ and $\dfrac14$.
\begin{enumerate}
\item Calculate the expectation $E(X)$. \hfill(1 pt)
\item Calculate $V(X)$ then the exact standard deviation $\sigma(X)$. \hfill(2 pts)
\item A game pays the gain $G=3X-4$. Calculate $E(G)$ and $V(G)$, then interpret the sign of $E(G)$. \hfill(2 pts)
\end{enumerate}
\mxneed{12}
\subsection*{Exercise 2 -- Comparing two models (5 points)}
Define on $\R$ the functions $f(x)=e^{2x}$ and $g(x)=1+2x$, then their difference $h=f-g$. The last question considers the values of $f$ at integer times, starting from index 0.
\begin{enumerate}
\item Calculate the derivative of $h=f-g$. \hfill(1 pt)
\item Study h and compare the curves. \hfill(2 pts)
\item For $u_n=e^{2n}$, give the common ratio and the sum of the first four terms. \hfill(2 pts)
\end{enumerate}
\mxneed{12}
\subsection*{Exercise 3 -- A displaced circle (4 points)}
In an orthonormal frame, consider the circle $\mathcal C$ with equation $(x-1)^2+(y+1)^2=4$ and the line $d$ with equation $y=-1$. Call $A$ and $B$ their points of intersection, $A$ being the one with the smaller x-coordinate.
\begin{enumerate}
\item Give the centre and radius. \hfill(1 pt)
\item Determine the intersections with the line. \hfill(1 pt)
\item For M on the circle, calculate $\vect{MA}\cdot\vect{MB}$ where A and B are these intersections; justify. \hfill(2 pts)
\end{enumerate}
